% !PRINTED 161
\weberpage{161}

This follows simply from the fact that a transposition, e.g. $(1, 2)$, corresponds to the substitution
\[
\begin{pmatrix} 0, & 1, & 0, & \dots & 0 \\ 1, & 0, & 0, & \dots & 0 \\ 0, & 0, & 1, & \dots & 0 \\ \cdot & \cdot & \cdot & \cdot & \cdot \\ 0, & 0, & 0, & \dots & 1 \end{pmatrix},
\]
whose determinant is $-1$, and from the fact that all permutations of the first kind can be composed from an even number of transpositions, and all permutations of the second kind from an odd number of transpositions.

\begin{center}
\S.\,40.

\gesperrt{The invariants of finite groups of linear substitutions.}
\end{center}

We now restrict our considerations to \gesperrt{finite groups} of linear homogeneous substitutions. The example of the permutation groups shows that such groups exist. Others we shall come to know later.

We denote such a group by $S$, and its substitutions by $A, B, C, \dots$ If in any function of the variables $x_1, x_2, \dots, x_n$ the variables $(x)$ are replaced by $A(x)$, we say that \gesperrt{the substitution $A$ is applied to the variables $(x)$}. We now first define the following concept:

1. \gesperrt{A form $\Phi(x_1, x_2, \dots, x_n)$ of $n$ variables is called an invariant or invariant form of the group $S$ if it remains unchanged when all the substitutions $A, B, C, \dots$ of the group $S$ are applied to the variables $(x)$.}

As before (Bd.\,I, \S.\,15), a form is here to be understood as an integral homogeneous function of the variables $x$.

We can also express this definition by the formulae:
\weq{1}{\Phi[A(x)] = \Phi[B(x)] = \Phi[C(x)] \dots}

That for every finite group $S$ there are invariant forms in arbitrary abundance is easy to see. One need only
%CONT>>
